% =============================================================================
%  3-Conclusiones.tex — Conclusiones y reflexión final
% =============================================================================

\section{Conclusiones}
\label{sec:conclusiones}

Se encontraron varios cambios durante este semestre que lejos de suponer una traba o un atraso, 
representaron una mejoría en los avances del mismo. El más significativo fue el cambio de 
paradigma de usar un enfoque 100\% urbano-arquitectónico a uno más 
computacional-matemático-urbano. De esta manera, se identifican 3 conclusiones clave:

\subsection{Conclusiones Nivel Documento}
\label{ssec:conclusiones-documento}

La implementación del sistema de composición \termino{LaTeX} para las fórmulas matemáticas
representó un avance contundente y acertado. Por el tipo de investigación, los casos de estudio y
las herramientas empleadas, fue la elección más adecuada: de haberse usado un procesador de textos
tradicional como Word o Google Docs, las referencias analizadas, las fórmulas matemáticas,
los diagramas y los bloques de código habrían representado un cuello de botella para el avance.

Por otro lado, se recurrió a herramientas de búsqueda profunda e inmersiva como \termino{NotebookLM}
y \termino{DeepSearch} para localizar artículos académicos citados en la bibliografía.

\subsection{Conclusiones Nivel Tecnológico}
\label{ssec:conclusiones-tecnologico}

La sección más difícil de implementar fue la programación de los algoritmos e indicadores
propuestos. Se emplearon los lenguajes de programación \termino{Python} y \termino{Go (Golang)}
para la codificación. El estado actual de los repositorios y el cambio de paradigma resultaron
ser un acierto.

\subsection{Conclusiones Nivel Maestría (Materias Cursadas)}
\label{ssec:conclusiones-nivel-maestría}

Aunque no pude completar mis estudios de manera presencial, las materias que más me ayudaron a
sustentar la investigación con bibliografía y casos de estudio fueron \termino{Sociología Urbana},
\termino{Desarrollo Urbano Sostenible} y \termino{Taller de Investigación: Megalópolis y Región}.
De esta última quiero destacar el gusto que sentí al poder citar cuatro de mis libros favoritos:
\begin{itemize}
    \item \textbf{Estado de crisis (Zygmunt Bauman):} Proporcionó la base sólida para comprender
            por qué el \termino{Estado de Bienestar} influyó en el abandono de la planeación urbana
            en la Ciudad de México entre las décadas de 1970 y 2000.
    \item \textbf{Las reglas del método sociológico (Émile Durkheim):} Base metodológica para analizar
            el deterioro del tejido social de la \zmvm.
    \item \textbf{Vida y muerte de las grandes ciudades (Jane Jacobs):} Funcionó como guía en tres ejes:
            \textit{Sociedad, Transporte y Espacios Públicos}, y sobre la importancia de no segregar a
            los ciudadanos ni reducirlos a categorías de menor jerarquía
            (\termino{Ciudadanos de primera, segunda y tercera clase}).
            Complementado con los informes de movilidad sustentable y desarrollo sostenible del Gobierno
            de la Ciudad de México e instituciones como la OMS y el INEGI.
    \item \textbf{Vigilar y castigar (Michel Foucault):} Base teórica sobre los peligros de la algoritmización
            y estandarización de la sociedad ante una realidad que supera el tejido urbano y económico, poniendo
            sobre la mesa la discusión de un \termino{estándar de personalidad} impuesto por las multinacionales
            y las necesidades geográficas. Complementado con \termino{La globalización} de Zygmunt Bauman.
\end{itemize}

\subsection{Conclusiones Nivel Personal}
\label{ssec:conclusiones-nivel-personal}

Aunque lejos de las aulas, pude disfrutar las lecturas que arriba mencioné, las cuales funcionaron
como un aliento en mi deteriorado estado de salud. Cuando podía leer el material que se me proporcionaba,
lograba distraer mi mente del malestar provocado por los medicamentos psiquiátricos. A la par, disfruté
relacionar las ideas de esos textos con otras de mis obras literarias favoritas:
\textit{La guía del autoestopista galáctico} (Douglas Adams),
\textit{¿Sueñan los androides con ovejas eléctricas?} (Philip K. Dick) y
\textit{La campana de cristal} (Sylvia Plath).

El cambio de paradigma fue como tener de nuevo las manos libres para hacer y materializar algo.
Aunque no me encontraba al 100\% de mi capacidad, programar las ideas matemáticas y plantear las
rutas de manera sistemática en lenguaje numérico resultó reconfortante.